%! AP Statistics — Unit 9 Cover Sheet
\documentclass[10pt]{article}
\usepackage{apstats-article}
\usepackage{apstats-boxes}

\begin{document}

% ── Full-bleed navy banner ────────────────────────────────────────────────────
\begin{tikzpicture}[remember picture, overlay]
  \fill[navy] (current page.north west)
    rectangle ([yshift=-1.4in]current page.north east);
\end{tikzpicture}
\vspace{-1.3in}

\begin{center}
  {\color{white}\bfseries\fontsize{24}{30}\selectfont AP Statistics}\\[5pt]
  {\color{goldacc}\bfseries\large Shepherd \quad 2026--2027}\\[10pt]
  {\color{white}\bfseries\Large Unit 9: Inference for Quantitative Data: Slopes}\\[3pt]
\end{center}

\vspace{0.18in}

% ── Unit overview ─────────────────────────────────────────────────────────────
\begin{tcolorbox}[colback=sky,colframe=navy,boxrule=0.9pt,arc=2mm,
    left=4mm,right=4mm,top=2mm,bottom=2mm]
  \small\textbf{Unit Overview.}\quad
  Unit 9 applies the inference framework to linear regression.
  Students construct and interpret confidence intervals for the slope of a
  population regression model, then perform significance tests to assess whether a
  linear relationship exists between two quantitative variables.
  The unit closes with a synthesis lesson in which students select the appropriate
  inference procedure across all procedures studied in the course.
\end{tcolorbox}

\vspace{0.14in}

% ── Lessons in this unit ──────────────────────────────────────────────────────
\renewcommand{\arraystretch}{1.35}
\begin{skillbox}[Lessons in This Unit]{goldbox}
\small
\begin{tabularx}{\linewidth}{@{} c >{\bfseries}p{0.46\linewidth} X @{}}
  \textbf{\#} & \textbf{Title} & \textbf{Focus} \\
  \hline
  9.1 & Introducing Statistics: Do Those Points Align?
      & Revisiting scatterplots and LSRL; sampling variability of slope; motivating inference. \\
  \hline
  9.2 & Confidence Intervals for the Slope of a Regression Model
      & $t$-interval for slope; standard error of $\hat{\beta}_1$; conditions; interpretation. \\
  \hline
  9.3 & Justifying a Claim About the Slope of a Regression Model Based on a Confidence Interval
      & Using a CI to support or refute a claim; four-step write-up. \\
  \hline
  9.4 & Setting Up a Test for the Slope of a Regression Model
      & Hypotheses; $t$-statistic for slope; conditions for inference on slope. \\
  \hline
  9.5 & Carrying Out a Test for the Slope of a Regression Model
      & Computing $t$ and $p$-value; decision; conclusion in context; full four-step. \\
  \hline
  9.6 & Skills Focus: Selecting an Appropriate Inference Procedure
      & Choosing among all confidence interval and significance test procedures; AP exam synthesis. \\
  \hline
\end{tabularx}
\end{skillbox}

\vspace{0.12in}

% ── Standards covered ─────────────────────────────────────────────────────────
\begin{skillbox}[AP Statistics Big Ideas \& Learning Objectives --- Unit 9]{greenbox}
\small
\begin{tabularx}{\linewidth}{@{} >{\bfseries}p{0.18\linewidth} X @{}}
  VAR-1.J   & Identify questions suggested by the relationship between two quantitative variables in context. \\
  VAR-7.A   & Describe a regression model for a population; identify parameters $\alpha$, $\beta_1$, and $\sigma$. \\
  UNC-4.AH  & Construct and interpret a confidence interval for the slope of a regression model. \\
  UNC-4.AI  & Justify a claim about a population slope based on a confidence interval. \\
  VAR-7.B; VAR-7.C  & State hypotheses for a test about slope; identify the $t$ test for slope and verify conditions. \\
  VAR-7.D; DAT-3.I  & Calculate the $t$ statistic and $p$-value for a test on slope; state a conclusion in context. \\
  DAT-3.J   & Select an appropriate inference procedure for a given situation involving quantitative data. \\
\end{tabularx}
\end{skillbox}

\end{document}
